\begin{statementbox}
\textbf{代数上的唯一性条件}\quad 给定成对样本 $(x_i,y_i)$ $(i=1,\ldots,n)$。只要
\[
 S_{xx}=\sum_{i=1}^n(x_i-\bar x)^2>0,
 \qquad \bar x=\frac1n\sum x_i,\quad\bar y=\frac1n\sum y_i.
\]
就存在唯一一条使纵向残差平方和最小的直线。

\textbf{建模上的适用性检查}\quad 散点整体大致呈线性趋势时，用直线描述才比较合理；这不是计算公式的前提。即使散点明显弯曲，只要 $S_{xx}>0$，仍可计算唯一的最小二乘直线，只是它可能遗漏原数据中的曲线结构。

% R038: 线性模型的三个前提情形（同尺度坐标）
\begin{center}
	\begin{tikzpicture}[x=1cm,y=1.25cm,>=Stealth,every node/.style={font=\small},axis/.style={->,black!70,thick},pt/.style={circle,fill=blue!75!black,inner sep=2.2pt}]
		\path[use as bounding box] (0,0) rectangle (12.4,4.8);
		\node[font=\large\bfseries,text=CStatement] at (6.2,4.45) {先看散点：趋势近似线性且 $x$ 必须有展开};
		\foreach \X/\title/\note in {0/{适合线性模型}/{散点围绕直线趋势},4.2/{直线会遗漏结构}/{明显弯曲},8.4/{斜率无法估计}/{$x$ 没有横向变化}} {
		\draw[rounded corners=5pt,fill=blue!2,draw=CStatement] (\X+.1,.25) rectangle (\X+3.9,3.9);
		\node[font=\bfseries,text=CStatement] at (\X+2,3.55) {\title};
		\draw[axis] (\X+.55,.75)--(\X+3.55,.75) node[right] {$x$};
		\draw[axis] (\X+.55,.75)--(\X+.55,3.05) node[above] {$y$};
		\node[align=center,text=black!70] at (\X+2, .42) {\note};
		}
		\draw[CStatement,very thick] (.75,1.05)--(3.35,2.85);
		\foreach \x/\y in {.9/1.15,1.35/1.45,1.8/1.65,2.25/2.25,2.75/2.42,3.15/2.9}{\node[pt] at (\x,\y) {};}
		\draw[orange!85!black,very thick] (4.95,1.05).. controls (6.1,1.15) and (6.65,3.0)..(7.65,3.0);
		\foreach \x/\y in {5.0/1.25,5.45/1.12,5.9/1.48,6.35/2.05,6.8/2.75,7.3/2.9}{\node[pt] at (\x,\y) {};}
		\draw[red!75!black,densely dashed,thick] (10.25,1.5)--(10.25,3.0);
		\foreach \y in {1.45,1.9,2.45,2.85}{\node[pt] at (10.25,\y) {};}
		\node[text=red!75!black,align=center] at (10.75,1.05) {$\sum(x_i-\bar x)^2=0$};
	\end{tikzpicture}
\end{center}

\textbf{最小二乘结论}\quad 在所有直线 $\hat y=a+bx$ 中，令纵向残差平方和
\[
 Q(a,b)=\sum_{i=1}^n\bigl(y_i-a-bx_i\bigr)^2=\sum_{i=1}^n e_i^2
\]
最小的直线为
\[
 \boxed{\hat b=\frac{\sum(x_i-\bar x)(y_i-\bar y)}{\sum(x_i-\bar x)^2}},
 \qquad \boxed{\hat a=\bar y-\hat b\bar x},\qquad \hat y=\hat a+\hat b x.
\]
这里 $e_i=y_i-\hat y_i$ 是沿 $y$ 轴方向的残差。特别地，回归线恒过样本中心 $(\bar x,\bar y)$。这里“唯一”指最优直线唯一；若 $S_{xx}=0$，仍有最优直线，但斜率不能由数据唯一确定。
\end{statementbox}
